% Supplementary Material: step-by-step derivations for
% "Diffusion anisotropy and entropy: reflections of each other".
% Every relation here is also checked numerically by analysis/identity_checks.py (check names in brackets).
\documentclass[11pt]{article}
\usepackage[margin=1in]{geometry}
\usepackage{amsmath,amssymb,mathtools}
\usepackage[hidelinks]{hyperref}
\usepackage{parskip}

\newcommand{\Div}{\mathrm{Div}}
\newcommand{\KL}{\Div_{\mathrm{KL}}}
\newcommand{\CV}{\mathrm{CV}}
\newcommand{\PR}{\mathrm{PR}}
\newcommand{\FA}{\mathrm{FA}}
\newcommand{\GFA}{\mathrm{GFA}}
\newcommand{\tr}{\mathrm{tr}}
\newcommand{\numcheck}[1]{\par{\small\raggedright\noindent\textit{Numerical check:} \texttt{#1}\par}}
\renewcommand{\theequation}{S\arabic{equation}}
\renewcommand{\thesection}{S\arabic{section}}

\title{Supplementary Material\\[4pt]
\large Step-by-step derivations for ``Diffusion anisotropy and entropy: reflections of each other''}
\author{Scott N. Hwang, Jonathon K. Maffie, Sangam G. Kanekar}
\date{}

\begin{document}
\maketitle

This document derives every relation stated in the paper, one step at a time, in the paper's notation.
Equation numbers without an S refer to the paper. Each derivation ends with the name of the check in
\texttt{analysis/identity\_checks.py} of the code repository that verifies it numerically.

\section{Setup}

Let $x_1,\ldots,x_K$ be positive numbers. Their mean is $\bar x = \frac1K\sum_i x_i$. The average of any
quantity over $i$ is written $\langle\cdot\rangle$, so $\bar x = \langle x\rangle$. The standard deviation uses
divisor $K$,
\begin{equation}
\sigma_x^2 = \tfrac1K\textstyle\sum_i (x_i-\bar x)^2 = \langle x^2\rangle - \bar x^2 ,
\label{eq:var}
\end{equation}
and the coefficient of variation is $\CV_x = \sigma_x/\bar x$. Dividing Eq.~\eqref{eq:var} by $\bar x^2$ gives
\begin{equation}
\frac{\langle x^2\rangle}{\bar x^2} = 1 + \CV_x^2 .
\label{eq:secondmoment}
\end{equation}

The proportions are $p_i = x_i/\sum_j x_j$. They are positive and sum to 1, so they form a probability
distribution. When all values are equal, every $p_i$ equals $1/K$. This is the uniform distribution $u$, with
$u_i = 1/K$. Because $\sum_j x_j = K\bar x$,
\begin{equation}
p_i = \frac{x_i}{K\bar x} .
\label{eq:p}
\end{equation}

The R\'enyi entropy of order $\alpha\ne1$ is $H_\alpha(p) = (1-\alpha)^{-1}\ln\sum_i p_i^{\alpha}$. The Shannon
entropy $H_1 = -\sum_i p_i\ln p_i$ is its limit as $\alpha\to1$ (Section~\ref{sec:order1}). The power mean of
order $r\ne0$ is $M_r(x) = \langle x^r\rangle^{1/r}$, and $M_0 = \exp\langle\ln x\rangle$ is the geometric mean.
$M_1 = \bar x$. The power mean increases with $r$, and all power means are equal only when all values are equal
(the power-mean inequality).

\section{The Jensen gap}

The power-mean Jensen gap of order $\alpha$ (Eq.~1) is
\begin{equation}
J_\alpha(x) = \frac{\langle x^{\alpha}\rangle}{\bar x^{\,\alpha}} - 1 .
\label{eq:J}
\end{equation}
Jensen's inequality says that for a convex function $\varphi$, $\langle\varphi(x)\rangle \ge \varphi(\bar x)$, with
equality only when all values are equal. For a concave function the inequality reverses. The function
$\varphi(x)=x^{\alpha}$ is convex for $\alpha>1$ and for $\alpha<0$, and concave for $0<\alpha<1$. So
$\langle x^\alpha\rangle - \bar x^{\,\alpha}$, and with it $J_\alpha$, is positive for $\alpha>1$ and $\alpha<0$ and
negative for $0<\alpha<1$, unless all values are equal. In both cases $J_\alpha$ has the sign of $\alpha(\alpha-1)$.
\numcheck{sign violations (gap sign, divergence sign by order)}

Since $\langle x^\alpha\rangle = M_\alpha^{\,\alpha}$ and $\bar x = M_1$, the gap is also
\begin{equation}
1 + J_\alpha = \Bigl(\frac{M_\alpha}{M_1}\Bigr)^{\alpha} .
\label{eq:Jpower}
\end{equation}

At order 2, Eq.~\eqref{eq:secondmoment} gives
\begin{equation}
J_2 = \frac{\langle x^2\rangle}{\bar x^2} - 1 = \CV_x^2 .
\label{eq:J2}
\end{equation}

\section{The identity, Eq.~2}
\label{sec:identity}

\textbf{Step 1, the power sum.} Raise Eq.~\eqref{eq:p} to the power $\alpha$ and sum over $i$:
\begin{equation}
\sum_i p_i^{\alpha} = \frac{\sum_i x_i^{\alpha}}{(K\bar x)^{\alpha}}
= \frac{K\langle x^{\alpha}\rangle}{K^{\alpha}\bar x^{\,\alpha}}
= K^{1-\alpha}\,\frac{\langle x^{\alpha}\rangle}{\bar x^{\,\alpha}}
= K^{1-\alpha}\,(1+J_\alpha) .
\label{eq:powersum}
\end{equation}

\textbf{Step 2, the entropy.} Take $(1-\alpha)^{-1}\ln$ of Eq.~\eqref{eq:powersum}:
\begin{equation}
H_\alpha = \frac{1}{1-\alpha}\Bigl[(1-\alpha)\ln K + \ln(1+J_\alpha)\Bigr] = \ln K + \frac{\ln(1+J_\alpha)}{1-\alpha} .
\end{equation}

\textbf{Step 3, the entropy deficit.} Move $\ln K$ to the left:
\begin{equation}
\ln K - H_\alpha = \frac{\ln(1+J_\alpha)}{\alpha-1} .
\label{eq:deficit}
\end{equation}

\textbf{Step 4, the divergence.} The R\'enyi divergence of $p$ from $u$ is
$\Div_\alpha(p\,\|\,u) = (\alpha-1)^{-1}\ln\sum_i p_i^{\alpha}u_i^{1-\alpha}$. Since $u_i^{1-\alpha} = K^{\alpha-1}$ for
every $i$,
\begin{equation}
\Div_\alpha(p\,\|\,u) = \frac{1}{\alpha-1}\Bigl[(\alpha-1)\ln K + \ln\sum_i p_i^{\alpha}\Bigr]
= \ln K - \frac{\ln\sum_i p_i^{\alpha}}{1-\alpha} = \ln K - H_\alpha .
\end{equation}

\textbf{Step 5, the power-mean form.} Insert Eq.~\eqref{eq:Jpower} into Eq.~\eqref{eq:deficit}:
\begin{equation}
\ln K - H_\alpha = \frac{1}{\alpha-1}\ln\Bigl(\frac{M_\alpha}{M_1}\Bigr)^{\alpha} = \frac{\alpha}{\alpha-1}\ln\frac{M_\alpha}{M_1} .
\end{equation}
Steps 3 to 5 together are Eq.~2.

\textbf{Step 6, the sign.} For $\alpha>1$, $M_\alpha\ge M_1$ and $\alpha/(\alpha-1)>0$, so the deficit is
nonnegative. For $0<\alpha<1$, $M_\alpha\le M_1$, so the logarithm is nonpositive, and $\alpha/(\alpha-1)<0$, so
the deficit is again nonnegative. In both cases it is zero only when all values are equal.

\textbf{Step 7, monotonicity.} For fixed $\alpha$, the right side of Eq.~\eqref{eq:deficit} is a strictly monotone
function of $J_\alpha$, increasing for $\alpha>1$ and decreasing for $\alpha<1$. Knowing one of the entropy deficit
and the Jensen gap of the same order therefore fixes the other.
\numcheck{Eq. 2 identity, every form, all orders | Eq. 2 derivation: sum p\^{}alpha = K\^{}(1-alpha)(1 + J\_alpha)}

\section{Order 1, the Shannon case}
\label{sec:order1}

\textbf{The limit.} At $\alpha=1$ both $\ln\sum_ip_i^{\alpha}$ and $1-\alpha$ vanish. By l'H\^opital's rule, with
$\frac{d}{d\alpha}\sum_i p_i^{\alpha} = \sum_i p_i^{\alpha}\ln p_i$ and $\sum_i p_i=1$,
\begin{equation}
H_1 = \lim_{\alpha\to1}\frac{\ln\sum_ip_i^{\alpha}}{1-\alpha} = \frac{\sum_i p_i\ln p_i}{-1} = -\sum_i p_i\ln p_i .
\end{equation}

\textbf{The deficit is the Kullback--Leibler divergence.} Since $\sum_ip_i=1$,
\begin{equation}
\ln K - H_1 = \sum_i p_i\ln K + \sum_i p_i\ln p_i = \sum_i p_i\ln\frac{p_i}{1/K} = \KL(p\,\|\,u) .
\end{equation}

\textbf{In terms of the values.} By Eq.~\eqref{eq:p}, $Kp_i = x_i/\bar x$, so
\begin{equation}
\ln K - H_1 = \sum_i p_i\ln\frac{x_i}{\bar x} = \sum_i p_i\ln x_i - \ln\bar x .
\label{eq:theil}
\end{equation}
This is Theil's index of the shares $p_i$.

\textbf{The Jensen gap of $x\ln x$.} For $\varphi(x)=x\ln x$, again using Eq.~\eqref{eq:p},
\begin{equation}
\frac{\langle x\ln x\rangle - \bar x\ln\bar x}{\bar x} = \sum_i\frac{x_i}{K\bar x}\ln x_i - \ln\bar x
= \sum_i p_i\ln x_i - \ln\bar x ,
\end{equation}
which is Eq.~\eqref{eq:theil}.

\textbf{The reverse direction.} With the roles of $p$ and $u$ exchanged,
\begin{equation}
\KL(u\,\|\,p) = \sum_i\frac1K\ln\frac{1/K}{p_i} = -\Bigl\langle\ln\frac{x}{\bar x}\Bigr\rangle = \ln\bar x - \langle\ln x\rangle = \ln\frac{M_1}{M_0} .
\label{eq:reverse}
\end{equation}
It can also be written $\KL(u\,\|\,p) = H(u,p) - \ln K$, where $H(u,p) = -\sum_i u_i\ln p_i$ is the cross-entropy
of $u$ relative to $p$, not an entropy of $p$.
\numcheck{uniform log gap ln(mean) - mean(ln) = KL(u||p) | Shannon deficit = Jensen gap of x ln x, scaled by mean}

\section{Other special orders}

\textbf{Order $\infty$.} As $\alpha\to\infty$, the largest proportion dominates $\sum_ip_i^{\alpha}$. If $m$
proportions share the largest value $p_{\max}$, then $\sum_ip_i^{\alpha} = m\,p_{\max}^{\alpha}(1+o(1))$, and
\begin{equation}
H_\infty = \lim_{\alpha\to\infty}\frac{\ln m + \alpha\ln p_{\max}}{1-\alpha} = -\ln p_{\max} ,
\qquad
\ln K - H_\infty = \ln(Kp_{\max}) = \ln\frac{x_{\max}}{\bar x} .
\end{equation}
\numcheck{order inf: ln K - H\_inf = ln(max x / mean x)}

\textbf{Order 2.} Eq.~\eqref{eq:powersum} with $\alpha=2$ and Eq.~\eqref{eq:J2} give
$\sum_ip_i^{2} = (1+\CV_x^2)/K$. Then
\begin{equation}
\ln K - H_2 = \ln(1+\CV_x^2), \qquad \PR = \frac{1}{\sum_ip_i^2} = \frac{K}{1+\CV_x^2},
\qquad \frac{\PR}{K} = \frac{1}{1+\CV_x^2} .
\label{eq:order2}
\end{equation}

\textbf{The chi-squared divergence.} Expanding the square and using $\sum_ip_i=1$,
\begin{equation}
\sum_i\frac{(p_i-u_i)^2}{u_i} = K\sum_i\Bigl(p_i^2 - \frac{2p_i}{K} + \frac{1}{K^2}\Bigr) = K\sum_ip_i^2 - 2 + 1
= (1+\CV_x^2) - 1 = \CV_x^2 .
\end{equation}
\numcheck{order 2: PR/K = 1/(1+CV\^{}2) | order 2: CV\^{}2 = chi\^{}2(p||u)}

\section{The geometric form, Eq.~3}

Let $\theta$ be the angle between the vector $x=(x_1,\ldots,x_K)$ and the isotropic direction $\mathbf 1=(1,\ldots,1)$.
Then
\begin{equation}
\cos^2\theta = \frac{(x\cdot\mathbf 1)^2}{|x|^2|\mathbf 1|^2} = \frac{(\sum_ix_i)^2}{K\sum_ix_i^2}
= \frac{K^2\bar x^2}{K\cdot K\langle x^2\rangle} = \frac{\bar x^2}{\langle x^2\rangle} = \frac{1}{1+\CV_x^2} = \frac{\PR}{K} ,
\end{equation}
using Eqs.~\eqref{eq:secondmoment} and \eqref{eq:order2}. Then $\tan^2\theta = 1/\cos^2\theta - 1 = \CV_x^2$, and
since the values are positive, $0\le\theta<\pi/2$ and $\tan\theta = \CV_x$. Also
$\sin^2\theta = 1-\cos^2\theta = \CV_x^2/(1+\CV_x^2)$.
\numcheck{PR/K = cos\^{}2 theta | CV = tan theta}

\section{Near isotropy, Eq.~4}
\label{sec:near}

Write $x_i = \bar x(1+\varepsilon_i)$. Averaging gives $\langle\varepsilon\rangle = 0$, and
$\langle\varepsilon^2\rangle = \sigma_x^2/\bar x^2 = \CV_x^2$. Expanding $(1+\varepsilon_i)^{\alpha}$ to second order,
\begin{equation}
1+J_\alpha = \langle(1+\varepsilon)^{\alpha}\rangle
= 1 + \alpha\langle\varepsilon\rangle + \tfrac12\alpha(\alpha-1)\langle\varepsilon^2\rangle + O(\varepsilon^3)
= 1 + \tfrac12\alpha(\alpha-1)\,\CV_x^2 + O(\varepsilon^3) .
\end{equation}
Since $\ln(1+J) = J + O(J^2)$, Eq.~\eqref{eq:deficit} gives
\begin{equation}
\ln K - H_\alpha = \frac{J_\alpha}{\alpha-1} + O(\varepsilon^3) = \frac{\alpha}{2}\,\CV_x^2 + O(\varepsilon^3) .
\end{equation}
Carrying the expansion of $(1+\varepsilon_i)^{\alpha}$ one term further adds
$\tfrac16\alpha(\alpha-1)(\alpha-2)\langle\varepsilon^3\rangle$ to $1+J_\alpha$. Since $J_\alpha^2$ is $O(\varepsilon^4)$,
\begin{equation}
\ln K - H_\alpha = \frac{\alpha}{2}\,\CV_x^2 + \frac{\alpha(\alpha-2)}{6}\,\langle\varepsilon^3\rangle + O(\varepsilon^4) .
\end{equation}
The cubic coefficient grows without bound with $\alpha$, so the approximation is not uniform in the order. At
$\alpha=\infty$ the deficit $\ln(x_{\max}/\bar x) = \ln(1+\varepsilon_{\max})\approx\varepsilon_{\max}$ is first order.
\numcheck{Eq. 4: J\_alpha / (alpha(alpha-1)/2 CV\^{}2) | Eq. 4: (ln K - H\_alpha) / (alpha/2 CV\^{}2) | Supplement S7: cubic term | Eq. 4: at alpha = infinity the deficit is first order}

\section{Divergences are not metrics}

A metric must be symmetric and satisfy the triangle inequality. At order $\tfrac12$,
$\Div_{1/2}(p\,\|\,q) = -2\ln\sum_i\sqrt{p_iq_i}$, which is symmetric in $p$ and $q$. It fails the triangle
inequality. For two outcomes, take $p=(0.9,0.1)$, $q=(0.5,0.5)$ and $r=(0.1,0.9)$. Then
$\Div_{1/2}(p\,\|\,q) = \Div_{1/2}(q\,\|\,r) = 0.2231$, but $\Div_{1/2}(p\,\|\,r) = -2\ln0.6 = 1.0217$, which exceeds
their sum. Other positive orders are not symmetric. For example, at order 1,
$\KL(p\,\|\,u) \ne \KL(u\,\|\,p)$ in general, as Section~\ref{sec:order1} shows.
\numcheck{order 1/2 is symmetric | orders 1, 2 and 4 are asymmetric | order 1/2 fails the triangle inequality | Supplement: triangle example}

\section{The diffusion tensor}

Take $x=\lambda$, the eigenvalues of the tensor, with $K=3$, mean $\bar\lambda$ and normalized eigenvalues
$p^{(\lambda)}_i = \lambda_i/\sum_j\lambda_j$.

\textbf{FA, Eq.~5.} Fractional anisotropy is $\FA^2 = \tfrac32\sum_i(\lambda_i-\bar\lambda)^2/\sum_i\lambda_i^2$. With
divisor 3, $\sum_i(\lambda_i-\bar\lambda)^2 = 3\sigma_\lambda^2 = 3\bar\lambda^2\CV_\lambda^2$ and, by
Eq.~\eqref{eq:secondmoment}, $\sum_i\lambda_i^2 = 3\langle\lambda^2\rangle = 3\bar\lambda^2(1+\CV_\lambda^2)$. So
\begin{equation}
\FA^2 = \frac32\,\frac{\CV_\lambda^2}{1+\CV_\lambda^2} .
\end{equation}
By Eq.~\eqref{eq:order2}, $\Div_2 = \ln3 - H_2 = \ln(1+\CV_\lambda^2)$, so $e^{-\Div_2} = 1/(1+\CV_\lambda^2)$ and
$1-e^{-\Div_2} = \CV_\lambda^2/(1+\CV_\lambda^2)$. Hence
\begin{equation}
\FA^2 = \tfrac32\bigl(1-e^{-\Div_2}\bigr) = \tfrac32\bigl(1-\PR_\lambda/3\bigr) = \tfrac32\sin^2\theta ,
\end{equation}
since $\PR_\lambda/3 = e^{-\Div_2}$ (Eq.~\eqref{eq:order2}) and $\sin^2\theta = \CV_\lambda^2/(1+\CV_\lambda^2)$.

\textbf{FA and the order-2 entropy.} From $\FA^2 = \tfrac32(1 - e^{H_2}/3)$, solve for $H_2$:
$e^{H_2} = 3 - 2\FA^2$, so
\begin{equation}
H_2 = \ln(3-2\FA^2), \qquad \ln3 - H_2 = -\ln\bigl(1 - \tfrac23\FA^2\bigr) .
\end{equation}
Near isotropy the second form is $\tfrac23\FA^2$ to leading order.
\numcheck{tensor orders, FA, QA and VR relations | ln 3 - H\_2(lambda) = -ln(1 - 2 FA\^{}2/3)}

\textbf{The order-1 member near isotropy.} Near isotropy $\FA^2 = \tfrac32\CV_\lambda^2/(1+\CV_\lambda^2) \approx
\tfrac32\CV_\lambda^2$, so $\CV_\lambda^2\approx\tfrac23\FA^2$. By Eq.~4 at $\alpha=1$ the order-1 member is then
$\approx\tfrac12\CV_\lambda^2 = \FA^2/3$ to leading order.
\numcheck{order-1 member / (FA\^{}2/3) near isotropy (FA < 0.05), median}

\textbf{The trace form of \"Ozarslan et al.} The trace-normalized tensor $R = \mathbf D/\tr\mathbf D$ has eigenvalues
$p^{(\lambda)}_i$, so $\tr R^2 = \sum_i(p^{(\lambda)}_i)^2 = e^{-H_2}$. Then
$\FA^2 = \tfrac32\bigl(1 - 1/(3\tr R^2)\bigr)$, which is $\FA = \sqrt{(3-1/\tr R^2)/2}$.
\numcheck{Ozarslan Eq. 11, FA from trace(R\^{}2), and trace(R\^{}2) = exp(-H\_2)}

\textbf{Procrustes anisotropy.} Dryden et al.\ define PA as FA computed on $q_i = \sqrt{\lambda_i}$. Since
$\sum_i(q_i-\bar q)^2 = \sum_iq_i^2 - 3\bar q^2$ and $\sum_iq_i^2 = \sum_i\lambda_i$,
\begin{equation}
\mathrm{PA}^2 = \frac32\Bigl(1 - \frac{(\sum_iq_i)^2}{3\sum_iq_i^2}\Bigr) .
\end{equation}
At order $\tfrac12$, with $u_i^{1/2} = 3^{-1/2}$,
$\Div_{1/2} = -2\ln\sum_i\sqrt{p^{(\lambda)}_i/3} = \ln3 - 2\ln\sum_i\sqrt{p^{(\lambda)}_i}$, so
$e^{-\Div_{1/2}} = (\sum_i\sqrt{p^{(\lambda)}_i})^2/3$. Since $\sqrt{p^{(\lambda)}_i} = q_i/\sqrt{\sum_j\lambda_j}$,
$(\sum_i\sqrt{p^{(\lambda)}_i})^2 = (\sum_iq_i)^2/\sum_iq_i^2$. Therefore
$\mathrm{PA}^2 = \tfrac32\bigl(1-e^{-\Div_{1/2}}\bigr)$.
\numcheck{Procrustes anisotropy (Dryden 2009): PA\^{}2 = (3/2)(1 - exp(-Div\_1/2))}

\textbf{Relative anisotropy.} In one common convention $\mathrm{RA} = \sqrt{\sum_i(\lambda_i-\bar\lambda)^2}/(\sqrt3\,\bar\lambda)
= \sigma_\lambda/\bar\lambda = \CV_\lambda$. Other conventions multiply this by a constant.

\textbf{Quantitative anisotropy, Eq.~6.} By Eq.~\eqref{eq:theil} with $x=\lambda$,
$\ln3 - H_1(p^{(\lambda)}) = \sum_ip^{(\lambda)}_i\ln\lambda_i - \ln\bar\lambda$. The von Neumann entropy of $R$ is
$-\tr(R\ln R) = -\sum_ip^{(\lambda)}_i\ln p^{(\lambda)}_i = H_1(p^{(\lambda)})$, because $R$ has eigenvalues
$p^{(\lambda)}_i$. So the order-1 member is $\ln3$ minus the von Neumann entropy of $R$.

\textbf{The volume ratio.} $\mathrm{VR} = \lambda_1\lambda_2\lambda_3/\bar\lambda^3$. The geometric mean is
$M_0 = (\lambda_1\lambda_2\lambda_3)^{1/3}$, so $\mathrm{VR} = (M_0/M_1)^3$, and by Eq.~\eqref{eq:reverse}
\begin{equation}
-\tfrac13\ln\mathrm{VR} = \ln\frac{M_1}{M_0} = \KL(u\,\|\,p^{(\lambda)}) .
\end{equation}

\textbf{FA and tensor mode fix the normalized eigenvalues.} The normalized eigenvalues sum to 1, so they have two
degrees of freedom. Write $d = p^{(\lambda)} - u$. Its components sum to zero, so $d$ lies in a plane, and $d$ is
orthogonal to $u$. Hence $\sum_i(p^{(\lambda)}_i)^2 = |u|^2 + |d|^2 = \tfrac13 + |d|^2$, and $|d|$ fixes
$\sum_i(p^{(\lambda)}_i)^2$ and therefore FA. In the plane, use the orthonormal basis
$e_1 = (2,-1,-1)/\sqrt6$, $e_2 = (0,1,-1)/\sqrt2$ and write $d = |d|(\cos\phi\,e_1 + \sin\phi\,e_2)$. The product of
the components of $d$ is $|d|^3\cos3\phi/(3\sqrt6)$, so the tensor mode, $3\sqrt6\prod_id_i/|d|^3$, equals
$\cos3\phi$. With $\lambda_1\ge\lambda_2\ge\lambda_3$, $\phi$ runs from 0 (linear, mode $+1$) to $\pi/3$ (planar, mode
$-1$), and $\cos3\phi$ is one-to-one on that range. So FA fixes $|d|$, mode fixes $\phi$, and together they fix
$p^{(\lambda)}$ and every entropy of it.
\numcheck{Supplement: mode = cos 3 phi and sum p\^{}2 = 1/3 + |d|\^{}2}

\textbf{The order-1 band at fixed FA.} Scale the eigenvalues so that $\lambda_1=1$, and let $0\le t\le1$. Linear
tensors are $(1,t,t)$, planar tensors $(1,1,t)$, and tensors with a zero eigenvalue $(1,t,0)$. By the definition of
FA in Eq.~5, the three families have
\begin{equation}
\FA^2 = \frac{(1-t)^2}{1+2t^2}, \qquad \FA^2 = \frac{(1-t)^2}{2+t^2}, \qquad \FA^2 = \frac{t^2-t+1}{1+t^2} .
\end{equation}
The planar family reaches its largest FA, $1/\sqrt2$, at $t=0$, the flat disc $(1,1,0)$. The zero-eigenvalue family
runs from the same disc at $t=1$ to the stick $(1,0,0)$, with $\FA=1$, at $t=0$. A fixed FA fixes
$\sum_i(p^{(\lambda)}_i)^2$, the index of coincidence. Harremo\"es and Tops{\o}e derived the joint range of the
Shannon entropy and the index of coincidence (cited in the paper). For three outcomes, its upper boundary is the
linear family. Its lower boundary is the planar family up to $\FA=1/\sqrt2$ and the zero-eigenvalue family above
it. We do not repeat their proof.
\numcheck{Supplement S9: FA of the linear, planar and zero-eigenvalue families | random eigenvalue triples inside the linear/planar/zero-eigenvalue interval | interval width at FA 0.2 | interval width at FA 0.5}

\section{The directional profile}

\textbf{GFA, Eq.~7.} For $N$ values $D_i$, Tuch's generalized fractional anisotropy is the standard deviation, with
divisor $N-1$, over the root mean square:
\begin{equation}
\GFA^2 = \frac{N\sum_i(D_i-\bar D)^2}{(N-1)\sum_iD_i^2} = \frac{N}{N-1}\,\frac{N\sigma_D^2}{N\langle D^2\rangle}
= \frac{N}{N-1}\,\frac{\CV_D^2}{1+\CV_D^2} .
\end{equation}
By Eq.~\eqref{eq:order2} and the geometric form this equals $\frac{N}{N-1}(1-\PR/N) = \frac{N}{N-1}\sin^2\theta$.
With $N=3$ the prefactor is $\tfrac32$ and this is Eq.~5.
\numcheck{Eq. 7: GFA\^{}2 (Tuch) = N/(N-1) CV\^{}2/(1+CV\^{}2) = N/(N-1)(1 - PR/N) = N/(N-1) sin\^{}2 theta}

\textbf{The CVD of Aja-Fern\'andez et al.} With $\log E_i = -bD_i$, their sample variance and second moment are
\begin{equation}
\mathcal V = \frac{N}{N-1}\bigl[\langle(\log E)^2\rangle - \langle\log E\rangle^2\bigr] = \frac{N}{N-1}\,b^2\sigma_D^2,
\qquad \langle(\log E)^2\rangle = b^2\langle D^2\rangle .
\end{equation}
So $\mathrm{CVD}^2 = \mathcal V/\langle(\log E)^2\rangle = \frac{N}{N-1}\,\sigma_D^2/\langle D^2\rangle = \GFA^2$.
\numcheck{CVD\^{}2 of Aja-Fernandez 2018 = GFA\^{}2 of the profile, Eq. (7)}

\textbf{The signal domain.} The signals $S_i/S_0 = e^{-bD_i}$ are positive, so the identity applies to them as
well. Near isotropy write $D_i = \bar D + \delta_i$ with $\langle\delta\rangle = 0$ and $\langle\delta^2\rangle =
\sigma_D^2$. Then $e^{-bD_i} = e^{-b\bar D}e^{-b\delta_i} \approx e^{-b\bar D}(1-b\delta_i)$. The signals have mean
$\approx e^{-b\bar D}$ and standard deviation $\approx e^{-b\bar D}b\sigma_D$, so their coefficient of variation is
$\approx b\sigma_D$. By Section~\ref{sec:near} the entropy deficits of the two domains are $\approx\frac{\alpha}{2}\CV_D^2$ and
$\approx\frac{\alpha}{2}(b\sigma_D)^2$, and their ratio is
\begin{equation}
\frac{(b\sigma_D)^2}{(\sigma_D/\bar D)^2} = (b\bar D)^2 .
\end{equation}
\numcheck{Sec 4.1 near isotropy: signal/diffusivity deficit ratio vs (b Dbar)\^{}2}

\section{The link to the tensor, Eq.~8}

For a single Gaussian compartment the profile is $D(\hat g) = \hat g^{\top}\mathbf D\,\hat g$. Averages over the unit
sphere of products of the components of $\hat g$ are
\begin{equation}
\langle g_ig_j\rangle = \tfrac13\delta_{ij}, \qquad
\langle g_ig_jg_kg_l\rangle = \tfrac{1}{15}\bigl(\delta_{ij}\delta_{kl} + \delta_{ik}\delta_{jl} + \delta_{il}\delta_{jk}\bigr) .
\end{equation}
The mean of the profile is
$\langle D\rangle = \sum_{ij}D_{ij}\langle g_ig_j\rangle = \tr\mathbf D/3 = \bar\lambda$. Its mean square is
\begin{equation}
\langle D^2\rangle = \sum_{ijkl}D_{ij}D_{kl}\langle g_ig_jg_kg_l\rangle
= \tfrac{1}{15}\bigl[(\tr\mathbf D)^2 + 2\tr\mathbf D^2\bigr]
= \tfrac{1}{15}\bigl[9\bar\lambda^2 + 2\textstyle\sum_i\lambda_i^2\bigr] .
\end{equation}
Its variance is
\begin{equation}
\langle D^2\rangle - \langle D\rangle^2 = \tfrac{1}{15}\bigl[2\textstyle\sum_i\lambda_i^2 - 6\bar\lambda^2\bigr]
= \tfrac{2}{15}\textstyle\sum_i(\lambda_i-\bar\lambda)^2 = \tfrac25\sigma_\lambda^2 ,
\end{equation}
using $\sum_i(\lambda_i-\bar\lambda)^2 = \sum_i\lambda_i^2 - 3\bar\lambda^2 = 3\sigma_\lambda^2$. Dividing by
$\bar\lambda^2$,
\begin{equation}
\CV_D^2 = \tfrac25\,\CV_\lambda^2 .
\end{equation}
Equal-weight averages over $N$ near-uniform directions approximate these sphere averages.

\textbf{The linear limit.} For $\lambda\propto(1,0,0)$, $\CV_\lambda^2 = 2$ and $\CV_D^2 = \tfrac45$. The normalized
second moments are $1+\CV_\lambda^2 = 3$ and $1+\CV_D^2 = \tfrac95$, whose ratio is $\tfrac35$, the value
\"Ozarslan et al.\ report for the profile against the eigenvalues.
\numcheck{Sec 4.2 sphere moments | CV\_D\^{}2 / CV\_lambda\^{}2 on a dense sphere, 50 random tensors | linear tensor: profile's normalized second moment / eigenvalues' = 3/5}

\section{The order-1 link between profile and tensor}

In the eigenvector basis, $D(\hat g) = \sum_iw_i(\hat g)\lambda_i$ with $w_i(\hat g) = (\hat g\cdot e_i)^2\ge0$ and
$\sum_iw_i = 1$. So each direction's diffusivity is a weighted mean of the eigenvalues. Let $f(x) = x\ln x$, which is
convex. Jensen's inequality for each direction gives $f(D(\hat g)) \le \sum_iw_i(\hat g)f(\lambda_i)$. Averaging over the
sphere, with $\langle w_i\rangle = \tfrac13$,
\begin{equation}
\langle f(D)\rangle \le \tfrac13\textstyle\sum_if(\lambda_i) = \langle f(\lambda)\rangle .
\end{equation}
The profile and the eigenvalues have the same mean $\bar\lambda$, so subtracting $f(\bar\lambda)$ and dividing by
$\bar\lambda$ turns both sides into order-1 deficits (Section~\ref{sec:order1}). The profile's order-1 deficit is
therefore at most the eigenvalues', with equality only at isotropy. Near isotropy both deficits are
$\approx\tfrac12\CV^2$, so their ratio tends to $\CV_D^2/\CV_\lambda^2 = \tfrac25$.
\numcheck{order-1 link: Jensen violations over the eigenvalue simplex | order-1 link: median ratio near isotropy (FA < 0.05)}

\section{The chain rule at order 1}

Split the $K$ values into groups $g$. Let $P_g = \sum_{i\in g}p_i$ be a group's proportion of the total and
$U_g = |g|/K$ its fraction of the values. Within a group, $p_{i|g} = p_i/P_g$ and $u_{i|g} = 1/|g|$, so
$p_i = P_gp_{i|g}$ and $u_i = U_gu_{i|g}$. Then
\begin{equation}
\KL(p\,\|\,u) = \sum_g\sum_{i\in g}P_gp_{i|g}\Bigl[\ln\frac{P_g}{U_g} + \ln\frac{p_{i|g}}{u_{i|g}}\Bigr]
= \KL(P\,\|\,U) + \sum_gP_g\,\KL(p_{\cdot|g}\,\|\,u_{\cdot|g}) ,
\end{equation}
using $\sum_{i\in g}p_{i|g} = 1$. The order-1 deficit is the proportion-weighted mean of the within-group deficits
plus the divergence of the groups' proportions from their fractions of the values.
\numcheck{order-1 deficit = share-weighted within-group deficits + between-group divergence}

\end{document}
